% Section IV: Method -- fair evaluation, features (with AP-affinity), models,
% tuning, the final-model rule and the metrics. Settings and grids:
% src/indoorloc/config.py; chosen settings: results/tuning/best_params.json.
% The example weights (0.27, 0.05, almost 0) and r_min = -105 dBm come from
% indoorloc.features.APAffinity fitted on the cleaned training data.
\section{Method}\label{sec:method}

\subsection{Fair Evaluation}\label{sec:fair}
Because a person records several scans in a row at one spot with one phone,
random CV almost always leaves training scans from the same spot, person and
phone for a held-out scan. We therefore use 5-fold \emph{user-group
cross-validation (CV)}~\cite{roberts2017crossvalidation}, the GroupKFold of
scikit-learn~\cite{pedregosa2011scikit} with the 16 user--device groups: each
fold holds out whole groups, so every model is tested on people and phones
that it has not seen, as for most of the test set (9 of its 11 phones are
new). Every choice (features, settings,
final model) is made with these folds on the training data only, and all that
a model learns from data, including the kept WAP columns and the AP-affinity
profiles, is learned again in each fold. The test set is used once, at the
very end: each model is fitted on all 18,719 training scans and predicts the
1,111 test scans. For comparison only, we also score the Baseline RF with
\emph{random CV} (shuffled 5-fold).

\subsection{Features}\label{sec:features}
We compare three feature sets: \emph{rssi}, the RSSI of the kept WAPs (464
columns on the full training data); \emph{rssi+stats}, which adds three
\emph{signal statistics} (the number of WAPs heard, the strongest reading and
the mean of the heard readings); and \emph{rssi+stats+affinity}, which adds our
eight \emph{AP-affinity} columns.

AP-affinity starts from a simple idea: if a phone hears a WAP that is mostly
heard in building~1 on floor~2, the phone is probably there too. It is learned
from the training rows only, in two steps. First, every WAP heard in at least
five training scans gets a \emph{profile}: the share of its scans in each
building (3 values), its mean floor and its mean position (2 values), each a
weighted mean over the training scans that hear it. A stronger reading usually
means that the phone was closer to the WAP, so it gets a larger weight
\begin{equation}
  w(r)=\left(\frac{r-r_{\min}}{-r_{\min}}\right)^{e},
  \label{eq:weight}
\end{equation}
where $r$ is the reading in dBm, $r_{\min}$ is 1~dB below the weakest training
reading ($-105$~dBm on the full training data), $e\approx2.718$, and $w=0$ if
the WAP is not heard. This ``powed'' formula~\cite{torres2015comprehensive}
gives $w=0.27$ at $-40$~dBm, $0.05$ at $-70$~dBm and almost 0 at $-100$~dBm.
Second, a scan $s$ is described by the weighted average of the profiles of the
known WAPs that it hears:
\begin{equation}
  \mathbf{p}_j=\frac{\sum_i w_{ij}\,\mathbf{t}_i}{\sum_i w_{ij}},\qquad
  \mathbf{a}(s)=\frac{\sum_j w_{sj}\,\mathbf{p}_j}{\sum_j w_{sj}},
  \label{eq:affinity}
\end{equation}
where $\mathbf{p}_j$ is the profile of WAP~$j$, $w_{ij}$ the weight of WAP~$j$
in scan~$i$, and $\mathbf{t}_i$ holds the targets of training scan~$i$ (three
0/1 building indicators, floor, $x$, $y$). The 6 values of $\mathbf{a}(s)$ and
the position of the strongest known WAP in the scan give 8 columns; a scan
that hears no known WAP gets the training means.

\emph{Missing values and scaling.} A WAP that is not heard gets $-110$~dBm in
both files, below every real reading. We add no separate heard/not-heard
column per WAP: for a tree, a split below $-105$~dBm already means ``not
heard'', and \texttt{n\_heard} counts the heard WAPs. We do not scale the RSSI:
trees use only the order of the values, and all RSSI columns share one unit
(dBm), so kNN treats them alike. Nothing is ever fitted on the test set.
Scaling, fitted on the training data only, is needed in Phase~2, since neural
networks train badly on unscaled inputs.

\subsection{Models}\label{sec:models}
Each method predicts building, floor and position with three separate models,
one per task: two classifiers and one regressor with two outputs ($x$, $y$),
built with scikit-learn~\cite{pedregosa2011scikit} and the XGBoost
library~\cite{chen2016xgboost}.
\begin{itemize}
\item \emph{Baseline RF}: a random forest~\cite{breiman2001random} with 300
  trees and default settings on rssi+stats.
\item \emph{kNN}~\cite{cover1967nearest} on rssi: the vote (building, floor)
  or mean position of the $k$ training scans with the closest RSSI values,
  plain or weighted by closeness.
\item \emph{Random Forest} on rssi+stats+affinity: 300 trees, each trained on
  a random sample of the scans with a random subset of the features tried at
  each split; the trees vote or average~\cite{breiman2001random}.
\item \emph{XGBoost} on rssi+stats+affinity: trees added one at a time, each
  correcting the errors of those before it~\cite{chen2016xgboost}; learning
  rate (0.1), row sample (0.85) and column sample (0.5) are fixed.
\end{itemize}
In formulas: kNN compares two scans $\mathbf{r},\mathbf{r}'$ (RSSI vectors)
with the Euclidean or the Manhattan distance,
\begin{equation}
  d_2=\sqrt{\textstyle\sum_j (r_j-r'_j)^2},\qquad
  d_1=\textstyle\sum_j |r_j-r'_j|.
  \label{eq:knn}
\end{equation}
A classification tree picks, at each node, the split that most lowers the Gini
impurity $G=1-\sum_c p_c^2$, where $p_c$ is the share of class $c$ in the node
($G=0$ when the node is pure); a regression tree lowers the squared error.
XGBoost builds its prediction step by step,
\begin{equation}
  F_m(\mathbf{x})=F_{m-1}(\mathbf{x})+\eta\,f_m(\mathbf{x}),
  \label{eq:boost}
\end{equation}
where each new tree $f_m$ is fitted to reduce the remaining loss and
$\eta=0.1$ is the learning rate.

\begin{table}[t]
  \caption{Tuned Settings: Values Tried and the Value Chosen by CV for Each
  Task (B: Building, F: Floor, P: Position)}
  \label{tab:grids}
  \centering
  \footnotesize
  \setlength{\tabcolsep}{4pt}
  \begin{tabular}{@{}llccc@{}}
    \toprule
    Setting & Values tried & B & F & P \\
    \midrule
    \multicolumn{5}{@{}l}{\emph{kNN} (28 settings)} \\
    n\_neighbors ($k$) & 1, 3, 5, 7, 9, 11, 15 & 1 & 5 & 3 \\
    weights & uniform, distance & unif. & unif. & dist. \\
    metric & Euclidean, Manhattan & Eucl. & Manh. & Eucl. \\
    \midrule
    \multicolumn{5}{@{}l}{\emph{Random Forest} (300 trees, 6 settings)} \\
    max\_features & sqrt, 0.2, 0.5 & sqrt & 0.5 & 0.5 \\
    min\_samples\_leaf & 1, 4 & 1 & 4 & 1 \\
    \midrule
    \multicolumn{5}{@{}l}{\emph{XGBoost} (6 settings)} \\
    max\_depth & 4, 6, 10 & 4 & 4 & 10 \\
    n\_estimators & 300, 600 & 300 & 300 & 600 \\
    \bottomrule
  \end{tabular}
\end{table}

\subsection{Tuning and the Final Model}\label{sec:tuning}
Each model has a small grid (Table~\ref{tab:grids}), and every setting is
scored on the same five folds. Because each task has its own model, each task
picks its own best setting: the highest CV accuracy for building and floor, the
lowest CV mean 2D error for position. Before using the test set, we fixed the
rule for the final model: the tuned model (kNN, Random Forest or XGBoost) with
the lowest CV mean 2D error.

\subsection{Metrics}\label{sec:metrics}
Building and floor accuracy are the shares of scans with the right building
and floor. For the position we report the mean absolute error (MAE) in $x$ and
$y$ and the mean and median \emph{2D error}, the straight-line distance between
the predicted and the true position: for test scan $i$ with true position
$(x_i,y_i)$ and prediction $(\hat x_i,\hat y_i)$,
\begin{equation}
  \mathrm{MAE}_x=\tfrac{1}{N}\textstyle\sum_i|\hat x_i-x_i|,\quad
  e_i=\sqrt{(\hat x_i-x_i)^2+(\hat y_i-y_i)^2}.
  \label{eq:metrics}
\end{equation}
To see where the error comes from, we
split the test scans into three cases (building wrong; building right but floor
wrong; both right) and give each case's share of scans, mean 2D error and share
of the total error. We also compare test scans from seen and new phones.
